\subsection{Infinite products of measures}

Let \((\mathbf{X}_{\lambda})_{\lambda \in \mathbf{L}}\) be a family of compact spaces, and for every \(\lambda \in \mathbf{L}\) let \(\mu_{\lambda}\) be a measure on \(\mathbf{X}_{\lambda}\) . We retain the notations of the Example of No.~5, so that in particular \(\mu_{\mathrm{J}} = \bigotimes_{\lambda \in \mathrm{J}} \mu_{\lambda}\) for every finite subset \(J\) of \(L\) .

PROPOSITION 9. --- Suppose that all of the measures \(\mu_{\lambda}\) are positive and that the family \((\mu_{\lambda}(1))_{\lambda \in L}\) is multipliable in \(\mathbf{R}_{+}\) (with product possibly 0). Then there exists one and only one measure \(\mu\) on X such that, for every finite subset J of L and every function \(f_{J} \in \mathcal{C}(X_J; C)\) ,

\[
\mu (f _ {\mathrm{J}} \circ \mathrm{pr} _ {\mathrm{J}}) = \mu_ {\mathrm{J}} (f _ {\mathrm{J}}) \prod_ {\lambda \in \mathrm{L} - \mathrm{J}} \mu_ {\lambda} (1).\tag{18}
\]

Moreover, the measure \(\mu\) is positive and its total mass is given by

\[
\mu (1) = \prod_ {\lambda \in L} \mu_ {\lambda} (1).\tag{19}
\]

Let \(\mathbf{F}\) be the vector space consisting of the functions on \(\mathbf{X}\) of the form \(f_{\mathrm{J}} \circ \mathrm{pr}_{\mathrm{J}}\) , where \(\mathbf{J}\) runs over the directed set of finite subsets of \(\mathbf{L}\) , and \(f_{\mathrm{J}} \in \mathcal{C}(\mathrm{X}_{\mathrm{J}}; \mathbf{C})\) ; we again say that \(\mathbf{F}\) is the space of continuous functions on \(\mathbf{X}\) that depend only on a finite number of variables. Lemma 3 of No.~5 shows that \(\mathbf{F}\) is dense in \(\mathcal{C}(\mathrm{X}; \mathbf{C})\) , which proves the uniqueness assertion. If \(\mu_{\lambda_0} = 0\) for some \(\lambda_0 \in \mathbf{L}\) then the measure \(\mu = 0\) meets the requirements, since in the second member of (18) we have \(\mu_{\mathrm{J}}(f_{\mathrm{J}}) = 0\) if \(\lambda_0 \in \mathbf{J}\) and \(\prod_{\lambda \in \mathbf{L} - \mathbf{J}} \mu_{\lambda}(1) = 0\) if \(\lambda_0 \notin \mathbf{J}\) . We can therefore suppose that \(\mu_{\lambda} \neq 0\) for all \(\lambda \in \mathbf{J}\) and, since the measures \(\mu_{\lambda}\) are positive, this implies that \(\mu_{\lambda}(1) \neq 0\) for all \(\lambda \in \mathbf{L}\) . Let us then set \(\mu_{\lambda}' = \mu_{\lambda}/\mu_{\lambda}(1)\) for every \(\lambda \in \mathbf{L}\) , so that \(\mu_{\lambda}'\) is a positive measure on \(X_{\lambda}\) such that \(\mu_{\lambda}'(1) = 1\) . It then follows from Prop. 8 of No.~5 that there exists a positive measure \(\mu'\) on \(X\) of total mass 1, such that \(\mu'(f_{\mathrm{J}} \circ \mathrm{pr}_{\mathrm{J}}) = \mu_{\mathrm{J}}'(f_{\mathrm{J}})\) for every finite subset \(J\) of \(L\) and every function \(f_{\mathrm{J}} \in \mathcal{C}(\mathrm{X}_{\mathrm{J}}; \mathbf{C})\) . The positive measure

\[
\mu = \left(\prod_ {\lambda \in L} \mu_ {\lambda} (1)\right) \mu^ {\prime}
\]

then meets the requirements, since

\[
\begin{array}{c} \mu_ {\mathrm{J}} (f _ {\mathrm{J}}) = \mu_ {\mathrm{J}} ^ {\prime} (f _ {\mathrm{J}}) \cdot \prod_ {\lambda \in \mathrm{J}} \mu_ {\lambda} (1), \\ \prod_ {\lambda \in \mathrm{L}} \mu_ {\lambda} (1) = \prod_ {\lambda \in \mathrm{J}} \mu_ {\lambda} (1) \cdot \prod_ {\lambda \in \mathrm{L} - \mathrm{J}} \mu_ {\lambda} (1). \end{array}
\]

The measure \(\mu\) defined in Prop. 9 is called the product measure of the family of positive measures \((\mu_{\lambda})_{\lambda \in L}\) and is denoted by \(\bigotimes_{\lambda \in L} \mu_{\lambda}\) .

COROLLARY. --- Assume that the conditions of Prop. 9 are verified, and let \((\mathrm{L}_{\rho})_{\rho\in\mathbb{R}}\) be a partition of L. Then each of the families of measures \((\mu_{\lambda})_{\lambda\in\mathrm{L}_{\rho}}\) admits a product measure, the family of product measures \(\left(\bigotimes_{\lambda\in\mathrm{L}_{\rho}}\mu_{\lambda}\right)_{\rho\in\mathbb{R}}\) admits a product measure, and

\[
\bigotimes_ {\rho \in \mathrm{R}} \left(\bigotimes_ {\lambda \in \mathrm{L} _ {\rho}} \mu_ {\lambda}\right) = \bigotimes_ {\lambda \in \mathrm{L}} \mu_ {\lambda}.\tag{20}
\]

This follows at once from the formulas (18) and (19) and the associativity of the product for multipliable families in \(R_{+}\) (GT, IV, §7, No.~5, Remark).

